% =============================================================================
\chapter{Experimental Results}
\label{ch:experiments}
% Colours of the direction-sweep maps (same palette as the images in data/).
\definecolor{sw0}{RGB}{68,1,84}
\definecolor{sw1}{RGB}{72,40,120}
\definecolor{sw2}{RGB}{62,74,137}
\definecolor{sw3}{RGB}{49,104,142}
\definecolor{sw4}{RGB}{38,130,142}
\definecolor{sw5}{RGB}{31,158,137}
\definecolor{sw6}{RGB}{53,183,121}
\definecolor{sw7}{RGB}{109,205,89}
\definecolor{sw8}{RGB}{253,231,37}
% Style of the exact zero-set slices in this chapter.
\pgfplotsset{zeroset/.style={axis equal image, empty line=jump, tick label style={font=\scriptsize},
    label style={font=\scriptsize}, title style={font=\small}, legend style={font=\scriptsize}, legend cell align=left}}
% =============================================================================
% PURPOSE: tests every prediction of Chapter 3 on all ten reference surfaces
% with the methods of Chapter 4, and measures how well Marching Cubes
% represents the neighbourhood of the singular points (the baseline for the
% comparison required by the assignment).

This chapter reports the experiments on the ten reference surfaces of
Table~\ref{tab:reference-surfaces}. Section~\ref{sec:exp-setup} lists the
perturbations. Sections~\ref{sec:exp-validation}--\ref{sec:exp-sweep} test
the predictions of Chapter~\ref{ch:theory}: convergence rates, singular
points, essential directions and discriminant directions.
Section~\ref{sec:exp-accuracy} measures how well the Marching Cubes meshes
represent the surfaces near their singular points, which is the baseline for
the comparison with other methods outlined in Section~\ref{sec:exp-comparison}.

% -----------------------------------------------------------------------------
\section{Setup}
\label{sec:exp-setup}
% PURPOSE: what was run, so that every table can be reproduced.

Every reference surface $F$ was combined with the same thirteen perturbations
\[
1,\ x,\ y,\ z,\ x+y,\ x+z,\ x+y+z,\ 2x-y+3z,\ x^2,\ x^2+y^2-z^2,\ xy,\ x^3,\ x^2y+y^2z+z^2x,
\]
covering the orders $0$ to $3$; with $G=F$, for which $S_k=S_0$ for all
$k\neq-1$ and every measured distance must vanish; and, for the four cubic cones,
with the balanced perturbation located by the direction sweep
(Section~\ref{sec:exp-sweep}). This gives $145$ families. Each family was
analysed at the resolutions $32$, $64$ and $128$ with the parameters of
Section~\ref{sec:method-parameters}: nine values of $|k|$ in
$[10^{-4},10^{-2}]$ on each side of zero, the local box $[-0.18,0.18]^3$ and
the domain $[-0.4,0.4]^3$. All results were produced by
\texttt{isosurface\_thesis\_experiments} in $17$ minutes on a ten-core laptop and are
stored as CSV files, one folder per surface, resolution and perturbation.

\begin{figure}[!h]
    \centering
    \begin{subfigure}{0.24\textwidth}\centering
        \thesisimage{pictures/ch5/surface_A1.png}{\textwidth}{3.2cm}{$x^2+y^2-z^2+xyz$, $k=0$}
        \caption{$A_1$}\end{subfigure}\hfill
    \begin{subfigure}{0.24\textwidth}\centering
        \thesisimage{pictures/ch5/surface_A2.png}{\textwidth}{3.2cm}{$x^2+y^2+z^3$, $k=0$}
        \caption{$A_2$}\end{subfigure}\hfill
    \begin{subfigure}{0.24\textwidth}\centering
        \thesisimage{pictures/ch5/surface_D4.png}{\textwidth}{3.2cm}{$x^2+y^3-3yz^2$, $k=0$}
        \caption{$D_4$}\end{subfigure}\hfill
    \begin{subfigure}{0.24\textwidth}\centering
        \thesisimage{pictures/ch5/surface_E6.png}{\textwidth}{3.2cm}{$x^3+y^3+z^3+xyz$, $k=0$}
        \caption{$\tilde E_6$}\end{subfigure}
    \caption[Representatives of the four types]{One reference surface of each type
    near its singular point (domain $[-0.4,0.4]^3$, resolution $128$).}
    \label{fig:types}
\end{figure}

% -----------------------------------------------------------------------------
\section{Validation of the Measurement}
\label{sec:exp-validation}
% PURPOSE: before the measurements can test the theory, the measuring
% instrument must be shown to be correct on cases with known answers, and
% the improvement over the mesh-to-mesh comparison must be quantified.

\paragraph{Known values.} For $G=F$ the measured local Hausdorff distance is
below $10^{-12}$ at every $k$ and resolution (``$=$'' in the tables). For the
Fermat cone and $G=1$ the distance from the origin to $S_k$ is exactly
$|k|^{1/3}$ (Section~\ref{sec:method-exact}); the measured values agree to
$10^{-10}$. The local Hausdorff distance itself behaves as predicted by
Proposition~\ref{prop:self-similar}: $d_a(k)=C|k|^{1/3}$ with $C=1.163$, the
Hausdorff distance between the model surface $\{F_+=\mp1\}$ and the cone. It is
larger than $1$ because the point of $S_k$ farthest from the cone lies in the
direction $(1,1,1)$, not on a coordinate axis. For the same reason the constant
depends on the perturbation and on the side: for $G=x$ it is $1.03$ for $k<0$
and $0.29$ for $k>0$.

\paragraph{Independence of the grid.} The exact distances do not depend on the
resolution of the mesh that supplied the starting points: across the resolutions
$32$, $64$ and $128$ the fitted exponents of the cubic cones differ by at most
$0.02$, and by less than $0.002$ for all linear perturbations. This is the decisive difference from a comparison of two meshes,
whose result depends on the grid (Figure~\ref{fig:mesh-vs-exact}). For
$F_++kx$ and $k>0$ the mesh-to-mesh comparison gives the exponents $0.38$,
$0.41$ and $0.44$ at the three resolutions, approaching the true value $0.5$
only slowly, because the deformation $0.29|k|^{1/2}$ is comparable to the grid
spacing over most of the range. The exact measurement gives $0.500$ at every
resolution. The asymmetry between $k<0$ and $k>0$ reported in earlier drafts
of this work was an artefact of the mesh comparison; the true rates are equal,
only the constants differ.

\begin{figure}[!h]
\centering
\begin{tikzpicture}
\begin{axis}[width=0.66\textwidth, height=0.46\textwidth, xmode=log, ymode=log,
    xlabel={$|k|$}, ylabel={$d_a(k)$}, legend pos=south east, legend cell align=left,
    legend style={font=\scriptsize}, grid=major, grid style={gray!25}, unbounded coords=discard]
\addplot[only marks, mark=*, mark size=1.4pt, black] table[x=absk, y=exact] {data/fermat_x_pos_res128.dat};
\addlegendentry{exact, every resolution}
\addplot[only marks, mark=o, mark size=1.6pt, red] table[x=absk, y=mesh] {data/fermat_x_pos_res32.dat};
\addlegendentry{mesh to mesh, $N=32$}
\addplot[only marks, mark=square, mark size=1.4pt, orange] table[x=absk, y=mesh] {data/fermat_x_pos_res64.dat};
\addlegendentry{mesh to mesh, $N=64$}
\addplot[only marks, mark=triangle, mark size=1.8pt, blue] table[x=absk, y=mesh] {data/fermat_x_pos_res128.dat};
\addlegendentry{mesh to mesh, $N=128$}
\addplot[black, dashed, domain=1e-4:1e-2, samples=2] {0.294*x^0.5};
\addlegendentry{$0.294\,|k|^{1/2}$}
\end{axis}
\end{tikzpicture}
\caption[Exact versus mesh-to-mesh distance]{Local Hausdorff distance for
$F_++kx$, $k>0$. The exact measurement follows $C|k|^{1/2}$ at all resolutions;
the distance between two meshes is dominated by their extraction errors when
the deformation is small and depends on the grid.}
\label{fig:mesh-vs-exact}
\end{figure}

% -----------------------------------------------------------------------------
\section{Convergence Rates}
\label{sec:exp-rates}
% PURPOSE: tests the scaling principle (3.5) on all surfaces and perturbations.

Table~\ref{tab:summary} summarizes all perturbations. Of the $268$ fitted
one-sided rates (all sides except $G=F$), $253$ have a $95\%$ confidence interval
that meets the predicted range~\eqref{eq:alpha-range}. Every stationary-point
count, $2592$ in total, is certified complete.

%%INSERT summary.tex

\paragraph{The cubic cones.} Table~\ref{tab:detail-fermat} lists the results
for the Fermat cone. All rates agree with $\alpha=1/(3-m)$ to within $0.002$: $\frac13$
for $G=1$, $\frac12$ for every linear $G$, and $1$ for all perturbations of
order two and three. The three other cones give the same exponents
($\frac13$, $\frac12$, $1$) with different constants; the exactness of the rates
is the self-similarity of Proposition~\ref{prop:self-similar}, which holds for
every homogeneous $F$ and $G$.

%%INSERT detail-fermat.tex

\paragraph{Quasi-homogeneous singularities.} For $A_1$, $A_2$ and $D_4$ the
weights differ between the variables, and the scaling principle predicts a
range of exponents (Tables~\ref{tab:detail-a1}--\ref{tab:detail-d4}). The data
show which end of the range is attained and why.
\begin{itemize}
    \item For the $A_1$ point, $G=1$ gives $\alpha=\frac12$ and every $G$ with
          $G(0)=0$ gives $\alpha=1$: such perturbations are inessential
          (their essential coordinates are zero), and their effect is a
          translation by $O(k)$ plus a constant of order $k^2$
          (Section~\ref{sec:vector}).
    \item For $x^2+y^2+z^3$ ($A_2$) and $G=z$, the predicted range is
          $[\frac12,\frac34]$, and both ends occur: $0.496$ for $k<0$ and $0.739$
          for $k>0$. For $k<0$ the polynomial $z^3+kz$ has the new real zeros
          $\pm\sqrt{-k}$, and the surface changes along the $z$-axis on the scale
          $|k|^{1/2}$; for $k>0$ there are no new zeros, and the change is
          visible only in the $x$ and $y$ directions, on the scale
          $\lambda^{1/2}=|k|^{3/4}$.
    \item For the same surface and $G=x$, the measured rate $0.66$ is the lower
          end $\frac23$. The perturbation is inessential, but it produces a thin
          spike, $(x+\frac k2)^2+y^2=\frac{k^2}{4}-z^3$, of height
          $(k^2/4)^{1/3}\propto|k|^{2/3}$ above the singular point
          (Figure~\ref{fig:spike-slice}).
\end{itemize}


\begin{figure}[!h]
\centering
\begin{tikzpicture}
\begin{axis}[zeroset, width=0.55\textwidth, xmin=-0.06, xmax=0.06, ymin=-0.06, ymax=0.06,
    xlabel={$x$}, ylabel={$z$}, legend pos=south east, xtick={-0.05,0,0.05}, ytick={-0.05,0,0.05}]
\addplot[black, thick] table[x=u, y=v] {data/slice_a2_spike_0.dat};
\addlegendentry{$k=0$}
\addplot[blue, thick] table[x=u, y=v] {data/slice_a2_spike_a.dat};
\addlegendentry{$k=0.003$}
\addplot[red, thick] table[x=u, y=v] {data/slice_a2_spike_b.dat};
\addlegendentry{$k=0.01$}
\addplot[gray, densely dashed, domain=-0.06:0.06, samples=2, forget plot] {0.0125};
\addplot[gray, densely dashed, domain=-0.06:0.06, samples=2, forget plot] {0.0292};
\addplot[gray, densely dotted, domain=-0.06:0.06, samples=2, forget plot] {0.00625};
\end{axis}
\end{tikzpicture}
\caption[The spike of an $A_2$ point]{Slice $y=0$ of $x^2+y^2+z^3+kx=0$. The
perturbation adds a spike above the singular point, of height
$(k^2/4)^{1/3}$ (dashed) and width about $k/2$. The dotted line marks one grid
spacing of the finest grid ($h=0.00625$): the spike is narrower than a cell
for every tested $k$, and Marching Cubes misses it.}
\label{fig:spike-slice}
\end{figure}

%%INSERT detail-a1.tex

%%INSERT detail-a2.tex

%%INSERT detail-d4.tex

\paragraph{Deviations.} The $15$ one-sided fits whose confidence interval
misses the prediction fall into two groups.
\begin{enumerate}
    \item \emph{Essential and inessential parts combined} ($10$ fits, $A_2$ and
          $D_4$, all but one for $k>0$). For example $x^2+y^2+z^3$ with $G=x+z$ gives $0.787$
          against the upper end $0.75$. The essential part $z$ deforms the
          surface on the scale $|k|^{3/4}$, while the inessential part $x$
          translates it by $O(|k|)$. The measured distance is the larger of the
          two contributions, and over the range $10^{-4}\le|k|\le10^{-2}$ the
          crossover between them bends the log-log line. The narrow confidence
          interval reflects the small scatter, not the curvature of the line
          (Section~\ref{sec:method-rates}); the prediction is asymptotic and
          is reached only for smaller $|k|$.
    \item \emph{Deformations below the resolution of the measurement} ($5$ fits,
          $D_4$ with $G=xy$, $x^2+y^2-z^2$ and $x^2y+y^2z+z^2x$). The deformation
          is of order $|k|\cdot10^{-4}$, and the measured distance levels off at
          about $1.5\cdot10^{-4}$. At a $D_4$ point the gradient vanishes to second
          order in two directions, the Lagrange system \eqref{eq:lagrange} is
          nearly singular, and the distance falls back to the spacing of the
          finest samples. These values are the resolution limit of the
          measurement, not a property of the surfaces.
\end{enumerate}

% -----------------------------------------------------------------------------
\section{Singular Points}
\label{sec:exp-singular}
% PURPOSE: tests Proposition 3.6 (linear perturbations of diagonal cones), the
% conservation of the Milnor number, and the behaviour at the singular point.

\paragraph{At the singular point.} In every experiment the behaviour of the
origin follows~\eqref{eq:at-origin}: it leaves the surface for $G=1$, becomes
regular when $G(0)=0\neq\nabla G(0)$, and stays singular when $G$ has order at
least two. For the Fermat cone and $G=x^2$ the certified search finds two
stationary points of multiplicity $4$, the origin and $(-\frac23k,0,0)$, both of
type $D_4$ (Hessian corank $2$), and only the origin on the zero level, as
derived in Section~\ref{sec:point-under-perturbation}.

\paragraph{New singular points.} Table~\ref{tab:singular-created} lists all
perturbations that create singular points away from the origin.
\begin{itemize}
    \item \emph{Linear perturbations of the Fermat cones} reproduce
          Proposition~\ref{prop:linear} exactly: $x+y$ and $x+z$ create two cusps
          ($A_2$, multiplicity $2$) for $k<0$ only, at distance $\sqrt{2|k|/3}$,
          i.e.\ with $\gamma=0.500$ (Figure~\ref{fig:cusps}); $x+y+z$ and $2x-y+3z$ create none; the
          balanced perturbations create two ordinary double points. For $F_-$,
          $x+z$ creates none, as the reflection~\eqref{eq:reflection} predicts.
    \item \emph{The quadratic cone $x^2+y^2-z^2$} added to the Fermat cones
          keeps the origin singular ($A_1$) and creates two further ordinary
          double points moving linearly, $\gamma=1$, for both signs of $k$; for
          the $D_4$ surface $x^2+y^3+z^3$ it creates one.
    \item \emph{An inessential perturbation can move a singular point.} For
          $x^2-y^2+z^3$ and $G=x+y$, $H_k=(x+y)(x-y+k)+z^3$ has its singular
          point at $(-\frac k2,\frac k2,0)$: the cusp is translated with rate
          $\gamma=1$, and the origin becomes regular.
    \item \emph{$D_4$ with a linear perturbation.} For $x^2+y^3-3yz^2$ and
          $G=y$, $k>0$, two ordinary double points appear with $\gamma=\frac12$:
          the cubic part $y^3-3yz^2$ has three real lines, and $+ky$ separates
          them (Figure~\ref{fig:d4-slice}).
\end{itemize}


\begin{figure}[!h]
\centering

\begin{tikzpicture}
\begin{axis}[zeroset, width=0.36\textwidth, xmin=-0.18, xmax=0.18, ymin=-0.18, ymax=0.18,
    xlabel={$x$}, ylabel={$y$}, title={$k=-0.01$}, xtick={-0.1,0,0.1}, ytick={-0.1,0,0.1}]
\addplot[blue, thick] table[x=u, y=v] {data/slice_fermat_xy_m.dat};
\addplot[only marks, mark=*, red, mark size=1.8pt] coordinates {(0.057735,-0.057735) (-0.057735,0.057735)};
\end{axis}
\end{tikzpicture}\hfill
\begin{tikzpicture}
\begin{axis}[zeroset, width=0.36\textwidth, xmin=-0.18, xmax=0.18, ymin=-0.18, ymax=0.18,
    xlabel={$x$}, title={$k=0$}, xtick={-0.1,0,0.1}, ytick={-0.1,0,0.1}, yticklabels={}]
\addplot[black, thick] table[x=u, y=v] {data/slice_fermat_xy_0.dat};
\addplot[only marks, mark=*, red, mark size=1.8pt] coordinates {(0,0)};
\end{axis}
\end{tikzpicture}\hfill
\begin{tikzpicture}
\begin{axis}[zeroset, width=0.36\textwidth, xmin=-0.18, xmax=0.18, ymin=-0.18, ymax=0.18,
    xlabel={$x$}, title={$k=+0.01$}, xtick={-0.1,0,0.1}, ytick={-0.1,0,0.1}, yticklabels={}]
\addplot[teal, thick] table[x=u, y=v] {data/slice_fermat_xy_p.dat};
\end{axis}
\end{tikzpicture}
\caption[Cusps of $F_++k(x+y)$]{The exact zero set of $F_++k(x+y)$ in the plane
$z=0$. For $k<0$ the cone opens at the origin and two cusps appear at
$(\pm s,\mp s,0)$, $s=\sqrt{-k/3}$ (red); for $k>0$ the surface is smooth.
The slices are computed from the polynomial, not from a mesh.}
\label{fig:cusps}
\end{figure}

%%INSERT singular-created.tex

\begin{figure}[!h]
\centering
\begin{tikzpicture}
\begin{axis}[zeroset, width=0.5\textwidth, xmin=-0.18, xmax=0.18, ymin=-0.18, ymax=0.18,
    xlabel={$y$}, ylabel={$z$}, legend pos=south east, xtick={-0.1,0,0.1}, ytick={-0.1,0,0.1}]
\addplot[black, thick] table[x=u, y=v] {data/slice_d4_y_0.dat};
\addlegendentry{$k=0$}
\addplot[blue, thick] table[x=u, y=v] {data/slice_d4_y_p.dat};
\addlegendentry{$k=0.004$}
\addplot[only marks, mark=*, red, mark size=1.8pt, forget plot] coordinates {(0,0.036515) (0,-0.036515)};
\end{axis}
\end{tikzpicture}
\caption[Two double points of a $D_4$ perturbation]{Slice $x=0$ of
$x^2+y^3-3yz^2+ky=0$. For $k=0$ the slice consists of the three lines
$y(y^2-3z^2)=0$ through the singular point. For $k>0$ the line $y=0$ and the
hyperbola $y^2-3z^2=-k$ remain, and they cross at the two ordinary double
points $(0,0,\pm\sqrt{k/3})$ (red), which converge to the origin with
$\gamma=\frac12$.}
\label{fig:d4-slice}
\end{figure}


In every case the multiplicities of all stationary points, real and complex,
add up to the Milnor number of $F$. The number of singular points never exceeds
it; the largest number observed is $3$ (origin plus two), against the bound $8$
for the cones.

\paragraph{A resolution limit of the count.} For $G=xy$ and the surfaces
$x^2+y^3+z^3$, $x^2+y^3-3yz^2$ and $x^3+y^3+z^2+x^2$ the table entry ``some $k$''
marks a stationary point that is reported as singular only for some values of
$k$. For $x^2+y^3+z^3+kxy$ it is the second stationary point
$(-\frac{k^3}{12},\frac{k^2}{6},0)$, at which $H_k=-k^6/432$. Its distance from
the origin, $k^2/6$, gives the measured rate $\gamma\approx2$; for
$|k|\lesssim0.008$ it is closer to the origin than the clustering radius
$10^{-5}$ and is counted together with the origin, and its value, below
$3\cdot10^{-15}$, is zero to within the tolerance of the search. Exactly, it is
not a singular point of $S_k$. The count of stationary points is still
complete; what the search cannot resolve is a point whose distance and value
are both below its tolerances.

% -----------------------------------------------------------------------------
\section{Essential Directions}
\label{sec:exp-essential}
% PURPOSE: tests "the perturbation as a vector" (3.6): the essential
% coordinates computed by the local algebra predict which perturbations only
% move the surface.

The second column of Tables~\ref{tab:detail-fermat}--\ref{tab:detail-d4}
gives the essential coordinates of each perturbation. They predict the
qualitative behaviour without any meshing.
\begin{itemize}
    \item A perturbation with zero essential part changes the surface only by
          a coordinate change at first order. Every such perturbation in the
          experiments has $\alpha=1$ (a translation or stretch of order $|k|$)
          or, at first order, nothing: $G=F$. The exceptions are the
          second-order effects of $A_1$ and $A_2$ (constant terms of order
          $k^2$), which give $\alpha=1$ for $A_1$ and $\alpha=\frac23$ for the
          spike of $A_2$.
    \item For the $A_2$ surface $x^3+y^3+z^2+x^2$ the essential part of
          $2x-y+3z$ is $-y$: only the cubic direction $y$ matters, the terms
          $2x+3z$ lie in the Tjurina ideal. Accordingly the measured exponents
          ($0.88$ for $k<0$, $0.53$ for $k>0$) are those of $G=y$ with the two
          sides exchanged ($0.74$ and $0.50$), shifted towards $1$ by the
          translation caused by $2x+3z$.
    \item For the cones all linear perturbations are essential, which is why
          they all give the rate $\frac12$; $x^2$ and the cubic perturbations are
          inessential and give $1$.
\end{itemize}

% -----------------------------------------------------------------------------
\section{Discriminant Directions}
\label{sec:exp-sweep}
% PURPOSE: the computed form of "G as a vector" for linear G: which directions
% c create singular points, compared with the exact formula for the cones.

Figure~\ref{fig:sweep} shows the direction sweep of linear perturbations at
$k=-0.01$ for one surface of each type. The colour is the number of real
stationary points near the origin; the black points are discriminant
directions located by bisection, at which the search confirms a singular point
of $S_k$. For $k=+0.01$ the maps are mirrored through the origin of the
sphere ($c\mapsto-c$).

\begin{figure}[!h]
\centering
\pgfplotsset{sweepmap/.style={width=0.5\textwidth, height=0.33\textwidth, xmin=-180, xmax=180, ymin=-1, ymax=1,
    enlargelimits=false, axis on top, xtick={-180,-90,0,90,180}, ytick={-1,0,1},
    xlabel={longitude of $c$ [deg]}, ylabel={$c_z$}, tick label style={font=\scriptsize},
    label style={font=\scriptsize}, title style={font=\small}}}
% The maps are images (one cell per grid direction, rendered from data/sweep_*.dat);
% the discriminant directions are drawn on top by pgfplots.
\begin{tikzpicture}
\begin{axis}[sweepmap, title={$A_1$: $x^2+y^2-z^2+xyz$}]
\addplot graphics[xmin=-180, xmax=180, ymin=-1, ymax=1] {data/sweep_A1.png};
\addplot[only marks, mark=*, mark size=0.8pt, black] table[x=lon, y=height] {data/sweepb_A1.dat};
\end{axis}
\end{tikzpicture}\hfill
\begin{tikzpicture}
\begin{axis}[sweepmap, title={$A_2$: $x^2-y^2+z^3$}]
\addplot graphics[xmin=-180, xmax=180, ymin=-1, ymax=1] {data/sweep_A2b.png};
\addplot[only marks, mark=*, mark size=0.8pt, black] table[x=lon, y=height] {data/sweepb_A2b.dat};
\end{axis}
\end{tikzpicture}

\begin{tikzpicture}
\begin{axis}[sweepmap, title={$D_4$: $x^2+y^3-3yz^2$}]
\addplot graphics[xmin=-180, xmax=180, ymin=-1, ymax=1] {data/sweep_D4b.png};
\addplot[only marks, mark=*, mark size=0.8pt, black] table[x=lon, y=height] {data/sweepb_D4b.dat};
\end{axis}
\end{tikzpicture}\hfill
\begin{tikzpicture}
\begin{axis}[sweepmap, title={$\tilde E_6$: $x^3+y^3+z^3$}]
\addplot graphics[xmin=-180, xmax=180, ymin=-1, ymax=1] {data/sweep_E6a.png};
\addplot[only marks, mark=*, mark size=0.8pt, black] table[x=lon, y=height] {data/sweepb_E6a.dat};
\end{axis}
\end{tikzpicture}

\smallskip
{\footnotesize Number of real stationary points:\enspace
\fcolorbox{black}{sw0}{\rule{0pt}{4pt}\rule{4pt}{0pt}}\,0\enspace
\fcolorbox{black}{sw1}{\rule{0pt}{4pt}\rule{4pt}{0pt}}\,1\enspace
\fcolorbox{black}{sw2}{\rule{0pt}{4pt}\rule{4pt}{0pt}}\,2\enspace
\fcolorbox{black}{sw3}{\rule{0pt}{4pt}\rule{4pt}{0pt}}\,3\enspace
\fcolorbox{black}{sw4}{\rule{0pt}{4pt}\rule{4pt}{0pt}}\,4\enspace
\fcolorbox{black}{sw5}{\rule{0pt}{4pt}\rule{4pt}{0pt}}\,5\enspace
\fcolorbox{black}{sw6}{\rule{0pt}{4pt}\rule{4pt}{0pt}}\,6\enspace
\fcolorbox{black}{sw7}{\rule{0pt}{4pt}\rule{4pt}{0pt}}\,7\enspace
\fcolorbox{black}{sw8}{\rule{0pt}{4pt}\rule{4pt}{0pt}}\,8}
\caption[Direction sweep]{Directions $c$ of linear perturbations $G=c\cdot p$
at $k=-0.01$, in the equal-area projection of the sphere. Colour: number of
real stationary points of $H_k$ near the origin. Black: discriminant
directions, for which $S_k$ has a singular point.}
\label{fig:sweep}
\end{figure}

For the Fermat cone the map has the structure predicted by
Proposition~\ref{prop:linear}. Real stationary points exist only in the octant
$c_x,c_y,c_z>0$ (eight of them) and on its faces and edges (fewer), and the
discriminant directions inside it are the curves
$\sum\varepsilon_i c_i^{3/2}=0$. On the meridian through $(1,1,0)$ the bisection
finds $c_z/c_x=1.587401052$, which is $2^{2/3}$ to the ten digits to which the
boundary is located; for $F_-$ it finds $-2^{2/3}$. For the cones without a
closed-form criterion it finds $x+y+2z$ for $x^3+y^3+z^3+xyz$ and
$x+y+2.658063392z$ for $x^3+2y^3+3z^3+2xyz$. Used as perturbations, all four
create two ordinary double points converging with $\gamma=\frac12$
(Table~\ref{tab:singular-created}): the predictions of the sweep are confirmed by
the independent family analysis.

For the $A_1$ point the discriminant directions form the cone
$c^{\top}A^{-1}c=0$ of Section~\ref{sec:vector} (perturbed at order $k$), a pair
of circles on the sphere; for the other types the map is more intricate, and,
unlike for the cones, it changes with $|k|$ (Proposition~\ref{prop:directions}
applies only to homogeneous $F$).

% -----------------------------------------------------------------------------
\section{Geometric Accuracy of Marching Cubes}
\label{sec:exp-accuracy}
% PURPOSE: the baseline for the comparison of methods (assignment): how far
% the Marching Cubes mesh is from the exact surface near the singular point,
% and what it misses.

Table~\ref{tab:extraction-error} shows the median extraction error, the
Hausdorff distance between the mesh and the exact surface in the box. It
decreases roughly in proportion to $h$ (order close to $1$) and stays below
half a grid cell for most families. The median hides the failures near the
singular points, which are the subject of this thesis.

%%INSERT extraction_error.tex

\paragraph{Missed features.} The largest errors occur where the surface has a
feature thinner than a grid cell. For $x^2+y^2+z^3$ and $G=x+y$ at
$|k|=0.0056$ the error is $0.025$ at resolution $128$, four grid cells: the spike
of Section~\ref{sec:exp-rates} is entirely missing from the mesh, and the
surface-to-mesh distance equals the height of the spike
(Figure~\ref{fig:spike}). The mesh-to-surface distance of the same mesh is only
$0.004$: every mesh point is close to the surface, but part of the surface is
far from the mesh. A measure that only checks mesh points cannot detect this.

\begin{figure}[!h]
\centering
\begin{tikzpicture}
\begin{axis}[width=0.62\textwidth, height=0.44\textwidth, xmode=log, ymode=log,
    xlabel={$|k|$, $k<0$}, ylabel={distance}, legend pos=north west, legend cell align=left,
    legend style={font=\scriptsize}, grid=major, grid style={gray!25}, unbounded coords=discard]
\addplot[black, mark=*, mark size=1.3pt] table[x=absk, y=exact] {data/a2_x_neg_res128.dat};
\addlegendentry{$d_a(S_k,S_0)$ (exact)}
\addplot[red, mark=o, mark size=1.5pt] table[x=absk, y=surftomesh] {data/a2_x_neg_res32.dat};
\addlegendentry{surface to mesh, $N=32$}
\addplot[blue, mark=triangle, mark size=1.7pt] table[x=absk, y=surftomesh] {data/a2_x_neg_res128.dat};
\addlegendentry{surface to mesh, $N=128$}
\addplot[black, dashed, domain=1e-4:1e-2, samples=2] {(x^2/4)^(1/3)};
\addlegendentry{spike height $(k^2/4)^{1/3}$}
\end{axis}
\end{tikzpicture}
\caption[A missed feature]{$x^2+y^2+z^3+kx$: the exact deformation is the
spike of height $(k^2/4)^{1/3}$. Where the surface-to-mesh distance coincides
with it, the mesh misses the spike completely.}
\label{fig:spike}
\end{figure}

\paragraph{Singular points.} At $k=0$ the singular point is the origin, a grid
node, and every mesh contains it (distance $0$). The moving singular points
are generally not grid nodes. Their distance to the mesh is up to $0.0106$ at
resolution $32$ (the balanced perturbation of the Fermat cone), about half a
cell, and decreases with $h$: $0.0064$ at $64$ and $0.0016$ at $128$. The mesh
passes near the singular point but does not contain it, and the triangles
around it describe a smooth surface, not a double point
(Section~\ref{sec:voxel-consequences}).

\begin{figure}[!h]
    \centering
    \begin{subfigure}{0.47\textwidth}\centering
        \thesisimage{pictures/ch5/balanced_mesh_32.png}{\textwidth}{4.5cm}{$F_+ + k(x+y+1.5874z)$, $k=-0.0032$, resolution 32, zoom on one of the two double points (mark the exact point)}
        \caption{resolution $32$}\end{subfigure}\hfill
    \begin{subfigure}{0.47\textwidth}\centering
        \thesisimage{pictures/ch5/balanced_mesh_128.png}{\textwidth}{4.5cm}{same family, $k$ and view, resolution 128}
        \caption{resolution $128$}\end{subfigure}
    \caption[Mesh near a moving singular point]{Marching Cubes near an ordinary
    double point that is not a grid node. The exact singular point is marked;
    the mesh does not contain it.}
    \label{fig:mesh-singular}
\end{figure}

% -----------------------------------------------------------------------------
\section{Self-Similarity}
\label{sec:exp-self-similarity}
% PURPOSE: a visual test of Proposition 3.2 (exact self-similarity).

For $F_++kx$ Proposition~\ref{prop:self-similar} gives $S_k=\sqrt{|k|}\,M_\sigma$
exactly. The zero sets for $k=0.01$, $0.04$ and $0.16$ must therefore be identical after
dividing the coordinates by $\lambda=\sqrt k$, which doubles from one value to
the next (Figure~\ref{fig:self-similar}). The same test can be made in the
viewer: extracting the three surfaces on the domains of half-size $0.4$, $0.8$
and $1.6$ with the same resolution gives identical meshes.

\begin{figure}[!h]
\centering
\begin{tikzpicture}
\begin{axis}[zeroset, width=0.55\textwidth, xmin=-4, xmax=4, ymin=-4, ymax=4,
    xlabel={$x/\sqrt{k}$}, ylabel={$y/\sqrt{k}$}, legend pos=south west]
\addplot[blue, line width=2.4pt, opacity=0.35] table[x=u, y=v] {data/slice_selfsim_a.dat};
\addlegendentry{$k=0.01$}
\addplot[red, line width=1.2pt] table[x=u, y=v] {data/slice_selfsim_b.dat};
\addlegendentry{$k=0.04$}
\addplot[black, densely dashed, thin] table[x=u, y=v] {data/slice_selfsim_c.dat};
\addlegendentry{$k=0.16$}
\end{axis}
\end{tikzpicture}
\caption[Self-similarity]{Slices $z=0$ of $F_++kx=0$ for three values of $k$,
drawn in the rescaled coordinates $p/\sqrt{k}$. The three curves coincide, as
Proposition~\ref{prop:self-similar} states: $S_k=\sqrt{k}\,M_+$.}
\label{fig:self-similar}
\end{figure}

% -----------------------------------------------------------------------------
\section{Limitations}
\label{sec:exp-limitations}
% PURPOSE: what the results do not show.

\begin{itemize}
    \item The exact distances are upper bounds that are attained when a start of
          the Lagrange iteration lies in the right basin. Near degenerate singular
          points ($D_4$) the iteration fails and the measurement has a floor of
          about $1.5\cdot10^{-4}$ in the box of half-size $0.18$.
    \item The rate fits cover two decades of $|k|$. Where two effects of
          different order combine, the fitted slope is not yet the asymptotic
          one, although its confidence interval is narrow.
    \item The certificate of completeness is numerical: it relies on distinct
          stationary points being farther apart than $10^{-5}$, which fails for
          the second stationary point of $x^2+y^3+z^3+kxy$.
    \item All singular points studied are at the origin, which is a grid node.
          A singular point between grid nodes is not represented by any mesh
          vertex, as the moving singular points show; this is the case the
          singularity-aware method has to handle.
    \item Only Marching Cubes has been evaluated so far.
\end{itemize}

% -----------------------------------------------------------------------------
\section{Comparison with Other Methods (Planned)}
\label{sec:exp-comparison}
% PURPOSE: placeholder for the comparison required by the assignment.

\begin{center}
\fbox{\parbox{0.92\textwidth}{\small\textbf{Work in progress.} The comparison
of Marching Cubes with Marching Tetrahedra, Dual Contouring and the method of
Section~\ref{sec:method-singularity-aware} has not been carried out yet. This
section lists what will be measured.}}
\end{center}

For every method and every family of Section~\ref{sec:exp-setup}, the
measurements of Section~\ref{sec:method-accuracy} will be repeated on the
method's mesh: the extraction error, the surface-to-mesh distance (missed
features), the distance from the singular points to the mesh, and the
topological checks of Section~\ref{sec:method-singularity-aware}. Since the
measurement pipeline accepts any triangle mesh, the results of
Table~\ref{tab:extraction-error} and Figure~\ref{fig:spike} are the Marching
Cubes column of that comparison.

% -----------------------------------------------------------------------------
\section{Conclusions}
\label{sec:exp-conclusions}

The experiments answer the questions of Chapter~\ref{ch:theory} for all ten
reference surfaces.
\begin{enumerate}
    \item \textbf{Rates.} The local Hausdorff distance between the exact
          surfaces decreases like $|k|^\alpha$ with $\alpha$ in the range
          predicted by the scaling principle in $253$ of $268$ cases; for
          homogeneous $F$ and $G$ the rate is exact, $\alpha=1/(d-m)$. The
          remaining cases are explained by combined essential and inessential
          parts, or by deformations below the resolution of the measurement.
    \item \textbf{Singular points.} The fate of the singular point is decided by
          $G(0)$ and $\nabla G(0)$; new singular points appear exactly where the
          theory predicts them, with the predicted type and rate, and their
          number never exceeds the Milnor number. All counts are certified.
    \item \textbf{Directions.} The essential coordinates of $G$ predict which
          perturbations only move the surface, and the direction sweep locates
          the perturbations that create singular points; the located balanced
          perturbations behave as predicted.
    \item \textbf{Extraction.} Marching Cubes represents the surfaces with an
          error of order $h$ away from the singular points, but it misses
          features thinner than a cell that start at stationary points, and it
          never contains singular points that are not grid nodes. These are the
          two failures that the singularity-aware method of
          Section~\ref{sec:method-singularity-aware} is designed to remove.
\end{enumerate}
